% ============================================================================
% Section: A/B Testing Framework
% ============================================================================
% Chapter 9: Production Engineering
% Topic: Rigorous validation framework for ensemble upgrades
% Author: Research Team
% Date: November 15, 2025
% Version: 1.0
%
% Purpose:
%   Document the comprehensive A/B testing infrastructure for validating
%   ensemble model improvements before production deployment.
%
% Key Topics:
%   - Hash-based deterministic routing
%   - Sequential testing with guardrails
%   - Statistical rigor (Diebold-Mariano, Bonferroni correction)
%   - Risk management and auto-stop mechanisms
%
% Dependencies:
%   - Section~\ref{sec:ensemble-evolution} (ensemble architecture)
%   - Section~\ref{sec:infrastructure-optimization} (latency requirements)
%   - Section~\ref{sec:temperature-scaling} (calibration methodology)
% ============================================================================

\section{A/B Testing Framework for Model Validation}
\label{sec:ab-testing-framework}

\subsection{Motivation: The Cost of Premature Deployment}

Production deployment of machine learning models for real-money betting requires exceptional rigor. Unlike academic research where false positives are publishable setbacks, premature deployment of an inferior model directly translates to financial loss. The challenge intensifies when evaluating incremental improvements: a proposed ensemble upgrade might show promising backtest results yet fail to generalize to live betting conditions due to overfitting, data drift, or implementation bugs.

\paragraph{Historical Context: The November 2025 XGBoost v4 Upgrade.}
In November 2025, we proposed upgrading ensemble v5.0 with XGBoost v4 Ridge Residuals (Section~\ref{sec:ensemble-evolution}), which demonstrated 88.31\% improvement over its XGBoost v2 baseline component on multi-season backtesting (2022--2024). However, backtest performance does not guarantee production performance:

\begin{itemize}[itemsep=3pt,parsep=0pt]
    \item \textbf{Implementation Risk:} Code differences between backtest and production environments
    \item \textbf{Data Drift:} 2025 season NFL dynamics may differ from 2022--2024 training data
    \item \textbf{Latency Degradation:} New model complexity could violate <200ms inference SLA
    \item \textbf{Calibration Breakdown:} Temperature scaling optimized for v5 may miscalibrate v6 predictions
    \item \textbf{Behavioral Changes:} Live betting markets react to our predictions, creating feedback loops
\end{itemize}

The cost of deploying an inferior model v6 over proven v5 was estimated at \$2,000--\$5,000 over 4 weeks (10\% ROI degradation on \$50,000 total bankroll). This motivated development of a rigorous A/B testing framework with statistical guardrails and early stopping mechanisms.

\subsection{A/B Test Design: Deterministic Hash-Based Routing}

Traditional A/B testing for web services uses randomized assignment at session-level, but sports betting predictions require \emph{deterministic assignment} to ensure reproducibility and audit compliance. Our framework assigns each (player, season, week) tuple to exactly one treatment group using cryptographic hashing.

\paragraph{Deterministic Assignment via MD5 Hashing.}
Given a prediction request for player $p$ in season $s$ and week $w$, we compute:

\begin{equation}
\begin{aligned}
\text{hash\_input} &= \text{concatenate}(p, s, w, \text{test\_name}) \\
\text{hash\_digest} &= \text{MD5}(\text{hash\_input}) \\
\text{hash\_value} &= \frac{\text{hash\_digest}[0:8]_{16} \bmod 10000}{10000} \in [0, 1) \\
\text{assignment} &= \begin{cases}
\text{treatment (B)} & \text{if } \text{hash\_value} < \rho \\
\text{control (A)} & \text{otherwise}
\end{cases}
\end{aligned}
\label{eq:ab-test-hash-assignment}
\end{equation}

where $\rho \in [0, 1]$ is the split ratio (typically $\rho = 0.5$ for 50/50 split), and $\text{hash\_digest}[0:8]_{16}$ extracts the first 8 hexadecimal characters of the MD5 digest as an integer.

\paragraph{Key Properties of Hash-Based Routing.}
This approach guarantees:

\begin{enumerate}[itemsep=2pt,parsep=0pt]
    \item \textbf{Determinism:} Same (player, season, week) always maps to same group
    \item \textbf{Uniformity:} Distribution approximates 50/50 split for large sample sizes
    \item \textbf{Reproducibility:} Assignment can be independently verified by auditors
    \item \textbf{No Cross-Contamination:} Each player-game receives exactly one prediction
    \item \textbf{Temporal Consistency:} Re-running predictions yields identical results
\end{enumerate}

The alternative---randomized assignment---would violate reproducibility (auditors couldn't verify which model produced which bet) and create cross-contamination risk (same player receiving both control and treatment predictions in same game).

\paragraph{Implementation Architecture.}
The \texttt{ABTestRouter} class (601 lines Python) implements this logic:

\begin{lstlisting}[language=Python, style=codestyle, caption={Deterministic A/B test assignment via MD5 hashing}]
import hashlib

class ABTestRouter:
    def _get_assignment(self, player_id: str, season: int,
                       week: int) -> tuple[str, str]:
        """Assign player-game to A/B group via MD5 hash."""
        # Create deterministic hash input
        hash_input = f"{player_id}_{season}_{week}_{self.test_name}"
        hash_digest = hashlib.md5(hash_input.encode()).hexdigest()

        # Convert to normalized value in [0, 1)
        hash_int = int(hash_digest[:8], 16)
        hash_value = (hash_int % 10000) / 10000.0

        # Assign to treatment if hash < split_ratio
        if hash_value < self.split_ratio:
            return ("B", "v6.0")  # Treatment
        else:
            return ("A", "v5.0")  # Control
\end{lstlisting}

This implementation routes predictions to ensemble v5.0 (control) or v6.0 (treatment) while logging all predictions to the \texttt{ml.predictions} table with \texttt{ab\_group} metadata for subsequent statistical analysis.

\subsection{Test Configuration: Ensemble v5.0 vs v6.0}

The November 2025 A/B test compared ensemble v5.0 (proven baseline) against v6.0 (XGBoost v4 Ridge upgrade). Table~\ref{tab:ab-test-config} summarizes the test configuration.

\begin{table}[htbp]
\centering
\small
\caption{A/B Test Configuration: Ensemble v5.0 (Control) vs v6.0 (Treatment)}
\label{tab:ab-test-config}
\begin{tabular}{@{}llS[table-format=2.2]S[table-format=2.1]@{}}
\toprule
\textbf{Group} & \textbf{Component} & {\textbf{Weight}} & {\textbf{Expected MAE (yds)}} \\
\midrule
\multirow{5}{*}{\textbf{Control (A)}} & \textit{Ensemble v5.0} & \multicolumn{2}{c}{\textit{(Proven Baseline)}} \\
\cmidrule(lr){2-4}
& State-Space v2 & 30.0 & {---} \\
& Hierarchical v2 & 25.0 & {---} \\
& XGBoost v2 Baseline & 20.0 & 40.20 \\
& Informative Priors v2 & 25.0 & {---} \\
\cmidrule(lr){2-4}
& \textit{Ensemble v5.0 Expected MAE} & \multicolumn{2}{c}{40.0 yards} \\
\midrule
\multirow{5}{*}{\textbf{Treatment (B)}} & \textit{Ensemble v6.0} & \multicolumn{2}{c}{\textit{(XGBoost v4 Ridge Upgrade)}} \\
\cmidrule(lr){2-4}
& State-Space v2 & 30.0 & {---} \\
& Hierarchical v2 & 25.0 & {---} \\
& XGBoost v4 Ridge & 20.0 & 4.70 \\
& Informative Priors v2 & 25.0 & {---} \\
\cmidrule(lr){2-4}
& \textit{Ensemble v6.0 Expected MAE} & \multicolumn{2}{c}{33.0 yards} \\
\midrule
\multicolumn{2}{l}{\textbf{Expected Improvement}} & \multicolumn{2}{c}{\textbf{7.0 yards (17.5\%)}} \\
\bottomrule
\end{tabular}
\begin{tablenotes}
\item \textbf{Test Duration:} 4 weeks (November 14 -- December 12, 2025), targeting 100+ games per group (2,000+ player-game predictions)
\item \textbf{Primary Hypothesis:} $H_1$: v6.0 reduces MAE by $\geq$5 yards vs v5.0 with statistical significance ($p < 0.05$)
\item \textbf{Split Ratio:} 50/50 deterministic assignment via MD5 hashing (Equation~\ref{eq:ab-test-hash-assignment})
\item \textbf{Key Difference:} XGBoost v4 Ridge (4.70 MAE on 2024 backtest) replaces XGBoost v2 Baseline (40.20 MAE), all other components identical
\item \textbf{Expected ROI Impact:} Control 5--7\%, Treatment 6--9\% (expected +1--2 pp improvement)
\end{tablenotes}
\end{table}

The test isolates the XGBoost v4 upgrade while holding all other ensemble components constant, enabling clean attribution of performance differences.

\subsection{Success Criteria: Multi-Metric Statistical Framework}

Deployment decisions require evidence beyond single-metric improvement. Our framework evaluates six metrics with explicit thresholds and statistical tests (Table~\ref{tab:ab-success-criteria}).

\begin{table}[htbp]
\centering
\small
\caption{A/B Test Success Criteria Matrix}
\label{tab:ab-success-criteria}
\begin{tabular}{@{}p{2.8cm}p{2.0cm}S[table-format=2.2]p{2.5cm}p{3.2cm}@{}}
\toprule
\textbf{Metric} & \textbf{Type} & {\textbf{Threshold}} & \textbf{Statistical Test} & \textbf{Rationale} \\
\midrule
\textbf{MAE} & Primary & {5.0} yards & Diebold-Mariano ($p < 0.05$) & Prediction accuracy improvement must be statistically significant and practically meaningful (5 yards = 12.5\% of 40-yard baseline) \\
\addlinespace[2pt]
\textbf{RMSE} & Secondary & {7.0} yards & Paired $t$-test ($p < 0.05$) & Validates improvement on outliers (RMSE penalizes large errors) \\
\addlinespace[2pt]
\textbf{Win Rate} & Business & {2.0} pp & Binomial proportion test ($p < 0.05$) & Directly impacts profitability (58\% → 60\% baseline → target) \\
\addlinespace[2pt]
\textbf{ROI} & Business & {1.0}\% & Bootstrap CI ($\alpha = 0.05$) & Primary business metric (5\% → 6\% target, must exceed Kelly threshold) \\
\addlinespace[2pt]
\textbf{Calibration} & Quality & {90.0}\% & Coverage test ($\pm$2 pp) & Uncertainty estimates must remain reliable (critical for Kelly bet sizing) \\
\addlinespace[2pt]
\textbf{Latency} & SLA & {200} ms & 95th percentile & Infrastructure constraint (cannot violate real-time betting SLA) \\
\bottomrule
\end{tabular}
\begin{tablenotes}
\item \textbf{Success Definition:} Treatment v6.0 must satisfy ALL primary criteria (MAE, Win Rate, ROI) plus NO degradation on secondary metrics (RMSE, Calibration, Latency)
\item \textbf{Failure Definition:} ANY primary criterion fails OR ANY guardrail violated (Section~\ref{sec:ab-guardrails}) → automatic rollback to control v5.0
\item \textbf{Inconclusive Definition:} $3 < \Delta\text{MAE} < 5$ yards OR $0.05 < p < 0.10$ → extend test by 2 weeks (50 additional games)
\item \textbf{Multiple Testing Correction:} Bonferroni correction applied to 6 metrics: $\alpha_{\text{corrected}} = 0.05 / 6 = 0.0083$ for family-wise error rate control
\item \textbf{Statistical Power:} $n = 100$ games per group provides 80\% power to detect 5-yard MAE improvement at $\alpha = 0.05$ (assumes $\sigma_{\text{diff}} = 12$ yards from historical data)
\end{tablenotes}
\end{table}

\paragraph{Diebold-Mariano Test for Forecast Accuracy.}
The primary statistical test is the Diebold-Mariano (DM) test~\citep{diebold1995comparing}, which evaluates whether two forecasting models have significantly different accuracy:

\begin{equation}
\begin{aligned}
d_t &= L(e_{t,\text{control}}) - L(e_{t,\text{treatment}}) \\
\bar{d} &= \frac{1}{n} \sum_{t=1}^{n} d_t \\
\text{DM} &= \frac{\bar{d}}{\sqrt{\widehat{\text{Var}}(\bar{d})}} \sim \mathcal{N}(0, 1) \quad \text{under } H_0
\end{aligned}
\label{eq:diebold-mariano}
\end{equation}

where $L(\cdot)$ is the loss function (absolute error for MAE), $e_{t,m}$ is the forecast error for model $m$ at time $t$, and $\widehat{\text{Var}}(\bar{d})$ uses Newey-West HAC standard errors to account for temporal autocorrelation in sports betting residuals.

The DM test null hypothesis $H_0: \mathbb{E}[d_t] = 0$ states that both models have equal expected loss. We reject $H_0$ (conclude treatment is superior) if DM statistic $< -1.96$ for $\alpha = 0.05$ two-tailed test (or $< -2.58$ after Bonferroni correction).

\subsection{Guardrails: Risk Management \& Auto-Stop Mechanisms}
\label{sec:ab-guardrails}

A/B tests in financial applications require aggressive guardrails to prevent catastrophic losses. Our framework implements three tiers of protection (Table~\ref{tab:ab-guardrails}).

\begin{table}[htbp]
\centering
\small
\caption{A/B Test Guardrails \& Auto-Stop Rules}
\label{tab:ab-guardrails}
\begin{tabular}{@{}p{2.5cm}S[table-format=2.1]lp{5.0cm}p{2.5cm}@{}}
\toprule
\textbf{Guardrail} & {\textbf{Threshold}} & \textbf{Severity} & \textbf{Trigger Condition} & \textbf{Auto-Stop Action} \\
\midrule
\textbf{MAE Degradation} & {10.0}\% & CRITICAL & Treatment MAE $>$ 44 yards (10\% worse than v5's 40 yards) & Immediate stop, rollback to v5, alert ops team \\
\addlinespace[2pt]
\textbf{ROI Degradation} & {-2.0}\% & CRITICAL & Treatment ROI $<$ 3\% (5\% baseline $-$ 2\%) for 2 consecutive days & Immediate stop, rollback to v5, alert ops team \\
\addlinespace[2pt]
\textbf{Win Rate Collapse} & {-5.0} pp & CRITICAL & Treatment win rate $<$ 53\% (58\% baseline $-$ 5\%) & Immediate stop, rollback to v5, alert ops team \\
\addlinespace[2pt]
\textbf{Latency SLA} & {200} ms & HIGH & 95th percentile latency $>$ 200ms for 6 consecutive hours & Alert ops team, consider load balancing, stop if $>$300ms \\
\addlinespace[2pt]
\textbf{Calibration Drift} & {5.0} pp & MEDIUM & Coverage drifts outside [85\%, 95\%] for 90\% nominal intervals & Alert data science, investigate temperature scaling, stop if $<$80\% \\
\addlinespace[2pt]
\textbf{Traffic Imbalance} & {10.0}\% & LOW & $|n_A - n_B| / (n_A + n_B) > 0.10$ (hash collision or bug) & Alert engineering, investigate routing logic \\
\bottomrule
\end{tabular}
\begin{tablenotes}
\item \textbf{Auto-Stop Mechanism:} All CRITICAL guardrails trigger automatic test termination and traffic rollback to control v5.0 within 60 seconds
\item \textbf{Monitoring Frequency:} Metrics refreshed every 1 hour, guardrails checked every 15 minutes, alerts via Grafana + PagerDuty
\item \textbf{Rollback SLA:} <2 minutes from guardrail violation to 100\% traffic on v5.0 (emergency kill switch bypasses gradual ramp-down)
\item \textbf{Historical Context:} Guardrails calibrated from 2022--2024 backtest volatility ($\sigma_{\text{MAE}} = 4.2$ yards, $\sigma_{\text{ROI}} = 1.5\%$)
\item \textbf{Financial Impact:} 10\% MAE degradation $\approx$ \$2,000 loss over 4 weeks (50\% of test traffic), ROI guardrail limits max loss to \$1,000
\end{tablenotes}
\end{table}

\paragraph{Implementation: Real-Time Guardrail Monitoring.}
Guardrails are implemented via Grafana alerting rules querying the \texttt{ml.predictions} table every 15 minutes:

\begin{lstlisting}[language=SQL, style=codestyle, caption={Real-time MAE degradation guardrail query}]
-- Alert if treatment MAE > 44 yards (10% degradation)
SELECT
  ab_group,
  AVG(ABS(prediction_mean - actual_value)) as mae
FROM ml.predictions
WHERE test_name = 'ensemble_v5_vs_v6_xgboost_ridge'
  AND actual_value IS NOT NULL
  AND predicted_at > NOW() - INTERVAL '2 hours'
GROUP BY ab_group
HAVING ab_group = 'B' AND mae > 44.0;
\end{lstlisting}

When triggered, Grafana alerts the operations team via PagerDuty and automatically executes the rollback script (\texttt{scripts/ab\_test\_emergency\_stop.sh}), which sets \texttt{split\_ratio = 0.0} to route 100\% traffic to control v5.0.

\subsection{Sequential Testing Strategy: Interim Analyses \& Early Stopping}

Fixed-horizon testing (wait 4 weeks, analyze once) wastes time when clear winner emerges early or when treatment underperforms. Our framework implements sequential testing with planned interim analyses (Table~\ref{tab:ab-sequential-schedule}).

\begin{table}[htbp]
\centering
\small
\caption{Sequential A/B Test Analysis Schedule}
\label{tab:ab-sequential-schedule}
\begin{tabular}{@{}clS[table-format=3.0]S[table-format=1.3]p{4.5cm}p{3.2cm}@{}}
\toprule
\textbf{Week} & \textbf{Analysis Type} & {\textbf{Min Games}} & {\textbf{$\alpha$-spending}} & \textbf{Decision Rule} & \textbf{Action} \\
\midrule
1 & Interim & 25 & {0.001} & Monitor only & Verify routing, check guardrails, no early stop \\
\addlinespace[2pt]
2 & Interim & 50 & {0.005} & Early stop if $p < 0.01$ AND $\Delta\text{MAE} > 7$ yards & If triggered: promote v6 to 100\%, else continue \\
\addlinespace[2pt]
3 & Interim & 75 & {0.010} & Monitor only & Check guardrails, estimate final power \\
\addlinespace[2pt]
4 & Final & 100 & {0.034} & Deploy if success criteria met & Final decision: deploy v6, keep v5, or extend test \\
\bottomrule
\end{tabular}
\begin{tablenotes}
\item \textbf{$\alpha$-Spending Function:} O'Brien-Fleming boundary~\citep{obrien1979multiple} with total $\alpha = 0.05$ distributed across 4 interim analyses to control Type I error inflation
\item \textbf{Cumulative $\alpha$:} Week 1 (0.001) + Week 2 (0.005) + Week 3 (0.010) + Week 4 (0.034) = 0.05 (family-wise error rate preserved)
\item \textbf{Early Stop Rationale:} Week 2 early stop requires stringent evidence ($p < 0.01$ vs final $p < 0.034$) to justify early deployment
\item \textbf{Power Monitoring:} Conditional power calculated at Week 3: if $<$ 50\% chance of success, recommend extending test by 2 weeks (50 games)
\item \textbf{Historical Context:} XGBoost v4 multi-season validation (2022--2024) showed consistent 4.70 $\pm$ 0.40 MAE, suggesting early stop likely at Week 2
\item \textbf{Type II Error Control:} $n = 100$ provides 80\% power to detect 5-yard improvement; early stop at $n = 50$ provides 60\% power (acceptable for $\Delta\text{MAE} > 7$ yards)
\end{tablenotes}
\end{table}

\paragraph{O'Brien-Fleming $\alpha$-Spending.}
Sequential testing requires careful control of family-wise Type I error rate (false positive rate) to avoid premature deployment due to random chance. The O'Brien-Fleming boundary~\citep{obrien1979multiple} allocates significance level $\alpha$ across $K$ interim analyses:

\begin{equation}
\alpha_k = 2 \left[ 1 - \Phi \left( \frac{z_{\alpha/2}}{\sqrt{k/K}} \right) \right], \quad k = 1, \ldots, K
\label{eq:obf-alpha-spending}
\end{equation}

where $\Phi(\cdot)$ is the standard normal CDF, $z_{\alpha/2} = 1.96$ for $\alpha = 0.05$ two-tailed test, and $k/K$ is the information fraction. For our 4-analysis schedule ($K = 4$):

\begin{align*}
\alpha_1 &= 2[1 - \Phi(1.96 / \sqrt{1/4})] = 0.001 \\
\alpha_2 &= 2[1 - \Phi(1.96 / \sqrt{2/4})] = 0.005 \\
\alpha_3 &= 2[1 - \Phi(1.96 / \sqrt{3/4})] = 0.010 \\
\alpha_4 &= 2[1 - \Phi(1.96 / \sqrt{4/4})] = 0.034
\end{align*}

This conservative early boundary ($\alpha_1 = 0.001$, $\alpha_2 = 0.005$) prevents spurious early stops while preserving most significance level for final analysis ($\alpha_4 = 0.034$).

\subsection{Implementation Architecture: Router, Logger, Analyzer}

The A/B testing framework consists of three integrated components (Figure~\ref{fig:ab-test-architecture}): router (prediction routing), logger (data collection), and analyzer (statistical testing).

\begin{figure}[htbp]
\centering
\begin{tikzpicture}[
    box/.style={rectangle, draw=black, thick, fill=white, minimum width=2.5cm, minimum height=1.0cm, align=center, font=\small},
    db/.style={cylinder, draw=black, thick, fill=white, minimum width=2.0cm, minimum height=1.2cm, align=center, font=\small, shape border rotate=90, aspect=0.25},
    arrow/.style={->, >=stealth, thick, color=blue!60!black},
    data/.style={->, >=stealth, thick, color=green!60!black, dashed},
    decision/.style={diamond, draw=black, thick, fill=yellow!20, minimum width=2.8cm, minimum height=1.5cm, align=center, font=\footnotesize, aspect=2},
    node distance=1.8cm
]

% Top: Prediction Request Flow
\node[box] (request) {Prediction\\Request};
\node[box, right=of request] (router) {\textbf{ABTestRouter}\\(Hash Assignment)};
\node[decision, right=of router] (assignment) {Hash $<$ 0.5?};
\node[box, above right=0.8cm and 1.5cm of assignment] (ensemble_v5) {Ensemble\\v5.0};
\node[box, below right=0.8cm and 1.5cm of assignment] (ensemble_v6) {Ensemble\\v6.0};
\node[box, right=3.5cm of assignment] (prediction) {Prediction\\Response};

% Database
\node[db, below=2.5cm of router] (db) {\texttt{ml.predictions}\\table};

% Bottom: Analysis Flow
\node[box, below=of db] (grafana) {Grafana\\Dashboards};
\node[box, right=of grafana] (analyzer) {\textbf{Statistical}\\Analyzer};
\node[decision, right=of analyzer] (decision) {Success\\Criteria\\Met?};
\node[box, above right=0.8cm and 1.0cm of decision] (deploy) {Deploy\\v6.0};
\node[box, below right=0.8cm and 1.0cm of decision] (rollback) {Keep\\v5.0};

% Arrows: Prediction flow
\draw[arrow] (request) -- node[above, font=\footnotesize] {$(p, s, w)$} (router);
\draw[arrow] (router) -- node[above, font=\footnotesize] {MD5 hash} (assignment);
\draw[arrow] (assignment) -- node[above right, font=\footnotesize, xshift=-5pt, yshift=5pt] {Yes: Group B} (ensemble_v6);
\draw[arrow] (assignment) -- node[above left, font=\footnotesize, xshift=5pt, yshift=5pt] {No: Group A} (ensemble_v5);
\draw[arrow] (ensemble_v5) -| (prediction);
\draw[arrow] (ensemble_v6) -| (prediction);

% Data logging
\draw[data] (router) -- node[right, font=\footnotesize] {Log} (db);
\draw[data] (ensemble_v5.south) -- ++(0,-0.8) -| node[near start, right, font=\footnotesize] {Store} (db.north);
\draw[data] (ensemble_v6.south) -- ++(0,-1.2) -| (db.north);

% Analysis flow
\draw[arrow] (db) -- node[above, font=\footnotesize] {Query} (grafana);
\draw[arrow] (grafana) -- node[above, font=\footnotesize] {Metrics} (analyzer);
\draw[arrow] (analyzer) -- node[above, font=\footnotesize] {DM test} (decision);
\draw[arrow] (decision) -- node[above right, font=\footnotesize, xshift=-5pt, yshift=3pt] {Yes ($p < 0.05$)} (deploy);
\draw[arrow] (decision) -- node[below right, font=\footnotesize, xshift=-5pt, yshift=-3pt] {No} (rollback);

% Guardrails
\node[box, fill=red!10, left=of grafana] (guardrails) {Guardrail\\Monitor};
\draw[arrow, color=red!60!black] (grafana) -- (guardrails);
\draw[arrow, color=red!60!black, dashed] (guardrails) |- node[near start, above, font=\footnotesize] {Auto-stop} (router);

\end{tikzpicture}
\caption{A/B Test Architecture: Hash-based routing with real-time monitoring and automated guardrails. Prediction requests are deterministically assigned to control (v5) or treatment (v6) via MD5 hashing. All predictions logged to database for statistical analysis. Grafana monitors guardrails and triggers auto-stop if degradation detected.}
\label{fig:ab-test-architecture}
\end{figure}

\paragraph{Component 1: ABTestRouter (Routing \& Assignment).}
The \texttt{ABTestRouter} class (601 lines Python) implements deterministic hash-based assignment (Equation~\ref{eq:ab-test-hash-assignment}) and routes predictions to appropriate ensemble:

\begin{itemize}[itemsep=2pt,parsep=0pt]
    \item \textbf{Input:} Player ID, season, week
    \item \textbf{Process:} Compute MD5 hash, assign to group A or B, route to ensemble v5 or v6
    \item \textbf{Output:} Prediction with metadata (\texttt{ab\_group}, \texttt{model\_version}, \texttt{test\_name})
    \item \textbf{Latency:} MD5 hashing adds 0.2ms overhead (negligible vs 167ms ensemble inference)
\end{itemize}

\paragraph{Component 2: Prediction Logger (Data Collection).}
All predictions are logged to the \texttt{ml.predictions} table with schema:

\begin{lstlisting}[language=SQL, style=codestyle, caption={A/B test prediction logging schema}]
CREATE TABLE ml.predictions (
  prediction_id SERIAL PRIMARY KEY,
  model_id INTEGER REFERENCES ml.model_registry(model_id),
  predicted_at TIMESTAMP NOT NULL,
  game_id TEXT NOT NULL,
  season INTEGER, week INTEGER,
  player_id TEXT, player_name TEXT,
  ab_group TEXT,          -- 'A' (control) or 'B' (treatment)
  model_version TEXT,     -- 'v5.0' or 'v6.0'
  test_name TEXT,         -- 'ensemble_v5_vs_v6_xgboost_ridge'
  prediction_mean NUMERIC,
  prediction_std NUMERIC,
  prediction_q10 NUMERIC, prediction_q90 NUMERIC,
  actual_value NUMERIC,   -- Filled post-game
  abs_prediction_error NUMERIC, -- |prediction - actual|
  calibrated BOOLEAN,     -- Within 90% interval?
  created_at TIMESTAMP DEFAULT NOW()
);
\end{lstlisting}

This schema enables efficient querying for A/B analysis, guardrail monitoring, and audit trail.

\paragraph{Component 3: Statistical Analyzer (Hypothesis Testing).}
The analyzer computes performance metrics and runs statistical tests:

\begin{lstlisting}[language=Python, style=codestyle, caption={Diebold-Mariano test implementation}]
import numpy as np
from scipy.stats import norm
from statsmodels.stats.sandwich_covariance import cov_hac

def diebold_mariano_test(errors_control, errors_treatment):
    """
    Test if treatment forecast accuracy differs from control.

    Returns:
        dm_stat: DM test statistic (reject H0 if < -1.96)
        p_value: One-tailed p-value for H1: treatment < control
    """
    # Compute loss differential (MAE)
    loss_diff = np.abs(errors_control) - np.abs(errors_treatment)

    # Mean differential
    mean_diff = np.mean(loss_diff)

    # HAC variance (Newey-West, lag=4 for weekly autocorrelation)
    var_diff = cov_hac(loss_diff.reshape(-1, 1), nlags=4)[0, 0]
    std_diff = np.sqrt(var_diff / len(loss_diff))

    # DM statistic
    dm_stat = mean_diff / std_diff
    p_value = norm.cdf(dm_stat)  # One-tailed (treatment < control)

    return dm_stat, p_value
\end{lstlisting}

The analyzer runs at each interim analysis (Table~\ref{tab:ab-sequential-schedule}), comparing DM $p$-value against O'Brien-Fleming boundary (Equation~\ref{eq:obf-alpha-spending}) to decide early stop or continuation.

\subsection{Sequential Testing Decision Flow}

Figure~\ref{fig:ab-sequential-flow} illustrates the week-by-week decision logic for the 4-week A/B test.

\begin{figure}[htbp]
\centering
\begin{tikzpicture}[
    box/.style={rectangle, draw=black, thick, fill=white, minimum width=3.0cm, minimum height=0.8cm, align=center, font=\footnotesize},
    decision/.style={diamond, draw=black, thick, fill=yellow!20, minimum width=3.5cm, minimum height=1.5cm, align=center, font=\footnotesize, aspect=2.2},
    outcome/.style={rectangle, draw=black, thick, fill=green!20, minimum width=2.5cm, minimum height=0.8cm, align=center, font=\footnotesize},
    stop/.style={rectangle, draw=black, thick, fill=red!20, minimum width=2.5cm, minimum height=0.8cm, align=center, font=\footnotesize},
    arrow/.style={->, >=stealth, thick, color=blue!60!black},
    node distance=1.5cm
]

% Week 1
\node[box] (week1) {\textbf{Week 1}\\$n = 25$ games};
\node[decision, below=of week1] (check1) {Guardrails\\OK?};
\node[stop, right=2.5cm of check1] (stop1) {STOP\\(Rollback v5)};
\node[box, below=of check1] (continue1) {Continue};

\draw[arrow] (week1) -- (check1);
\draw[arrow] (check1) -- node[above, font=\footnotesize] {No} (stop1);
\draw[arrow] (check1) -- node[right, font=\footnotesize] {Yes} (continue1);

% Week 2
\node[box, below=of continue1] (week2) {\textbf{Week 2}\\$n = 50$ games};
\node[decision, below=of week2] (check2) {$p < 0.01$\\AND\\$\Delta\text{MAE} > 7$?};
\node[outcome, right=2.5cm of check2] (deploy2) {DEPLOY v6\\(Early Win)};
\node[decision, below=of check2] (guard2) {Guardrails\\OK?};
\node[stop, right=2.5cm of guard2] (stop2) {STOP\\(Rollback v5)};
\node[box, below=of guard2] (continue2) {Continue};

\draw[arrow] (continue1) -- (week2);
\draw[arrow] (week2) -- (check2);
\draw[arrow] (check2) -- node[above, font=\footnotesize] {Yes} (deploy2);
\draw[arrow] (check2) -- node[right, font=\footnotesize] {No} (guard2);
\draw[arrow] (guard2) -- node[above, font=\footnotesize] {No} (stop2);
\draw[arrow] (guard2) -- node[right, font=\footnotesize] {Yes} (continue2);

% Week 3
\node[box, below=of continue2] (week3) {\textbf{Week 3}\\$n = 75$ games};
\node[decision, below=of week3] (check3) {Guardrails\\OK?};
\node[stop, right=2.5cm of check3] (stop3) {STOP\\(Rollback v5)};
\node[box, below=of check3] (continue3) {Continue};

\draw[arrow] (continue2) -- (week3);
\draw[arrow] (week3) -- (check3);
\draw[arrow] (check3) -- node[above, font=\footnotesize] {No} (stop3);
\draw[arrow] (check3) -- node[right, font=\footnotesize] {Yes} (continue3);

% Week 4 (Final)
\node[box, below=of continue3] (week4) {\textbf{Week 4 (Final)}\\$n = 100$ games};
\node[decision, below=of week4] (final) {Success\\Criteria\\Met?};
\node[outcome, right=2.5cm of final] (deploy4) {DEPLOY v6\\(Success)};
\node[stop, below=of final] (keep5) {KEEP v5\\(Failure)};

\draw[arrow] (continue3) -- (week4);
\draw[arrow] (week4) -- (final);
\draw[arrow] (final) -- node[above, font=\footnotesize] {Yes} (deploy4);
\draw[arrow] (final) -- node[right, font=\footnotesize] {No} (keep5);

\end{tikzpicture}
\caption{Sequential A/B test decision flow over 4 weeks. Guardrails checked every week; early stop at Week 2 if overwhelming evidence ($p < 0.01$, $\Delta\text{MAE} > 7$ yards); final decision at Week 4 based on success criteria (Table~\ref{tab:ab-success-criteria}).}
\label{fig:ab-sequential-flow}
\end{figure}

\paragraph{Week 1: Guardrail Validation.}
First week focuses on system validation:
\begin{itemize}[itemsep=2pt,parsep=0pt]
    \item Verify 50/50 traffic split (hash routing correctness)
    \item Confirm database logging working (no dropped predictions)
    \item Validate latency <200ms (no SLA violations)
    \item Check calibration stability (temperature scaling applies correctly)
    \item No statistical testing ($\alpha_1 = 0.001$ too stringent for $n = 25$)
\end{itemize}

\paragraph{Week 2: Early Stop Opportunity.}
If treatment shows overwhelming superiority ($p < 0.01$ AND $\Delta\text{MAE} > 7$ yards), deploy early:
\begin{itemize}[itemsep=2pt,parsep=0pt]
    \item Requires stringent evidence ($p < 0.01$ vs final threshold $p < 0.034$)
    \item Effect size must exceed expected improvement (7 yards > 5 yard threshold)
    \item Avoids 2 additional weeks of conservative baseline bets
    \item Risk: Conditional power at $n = 50$ is 60\% (acceptable for large effect sizes)
\end{itemize}

\paragraph{Week 3: Conditional Power Analysis.}
Estimate probability of success at final analysis:
\begin{equation}
\text{CP} = \Phi \left( \frac{\bar{d}_3 \sqrt{n_4}}{s_3} - z_{\alpha_4} \right)
\label{eq:conditional-power}
\end{equation}
where $\bar{d}_3$ is observed mean difference at $n_3 = 75$, $s_3$ is standard deviation, $n_4 = 100$ is final sample size, and $z_{\alpha_4} = 1.82$ for $\alpha_4 = 0.034$.

If CP $< 50\%$ (unlikely to succeed), recommend extending test by 2 weeks (50 additional games).

\paragraph{Week 4: Final Decision.}
Compare all 6 metrics (Table~\ref{tab:ab-success-criteria}) against thresholds:
\begin{itemize}[itemsep=2pt,parsep=0pt]
    \item \textbf{Deploy v6:} All primary criteria met (MAE, Win Rate, ROI) AND no degradation on secondary metrics
    \item \textbf{Keep v5:} Any primary criterion fails OR any guardrail violated
    \item \textbf{Extend Test:} Inconclusive ($3 < \Delta\text{MAE} < 5$ yards OR $0.05 < p < 0.10$)
\end{itemize}

\subsection{Monitoring Infrastructure: Real-Time Dashboards}

Grafana dashboards provide real-time visibility into A/B test performance (Table~\ref{tab:ab-grafana-panels}).

\begin{table}[htbp]
\centering
\small
\caption{Grafana A/B Test Monitoring Panels}
\label{tab:ab-grafana-panels}
\begin{tabular}{@{}p{3.5cm}p{5.5cm}p{4.0cm}@{}}
\toprule
\textbf{Panel Name} & \textbf{Query} & \textbf{Alert Trigger} \\
\midrule
\textbf{MAE by Group} & \texttt{AVG(ABS(prediction - actual))} grouped by \texttt{ab\_group} over rolling 24 hours & Treatment MAE $>$ 44 yards (10\% degradation) \\
\addlinespace[2pt]
\textbf{Cumulative Predictions} & \texttt{COUNT(*)} by \texttt{ab\_group}, partitioned by date & Traffic imbalance $>$ 10\% (hash routing bug) \\
\addlinespace[2pt]
\textbf{Win Rate Comparison} & Fraction of correct over/under predictions by group & Treatment win rate $<$ 53\% (5pp collapse) \\
\addlinespace[2pt]
\textbf{ROI Time Series} & Kelly-weighted returns by group, cumulative over test duration & Treatment ROI $<$ 3\% for 2 consecutive days \\
\addlinespace[2pt]
\textbf{Calibration Coverage} & Fraction of actuals within 90\% prediction intervals by group & Coverage $<$ 85\% or $>$ 95\% for 6 hours \\
\addlinespace[2pt]
\textbf{Inference Latency (p95)} & 95th percentile latency from \texttt{predicted\_at} timestamp by group & Latency $>$ 200ms for 6 consecutive hours \\
\bottomrule
\end{tabular}
\begin{tablenotes}
\item \textbf{Refresh Interval:} All panels refresh every 5 minutes, guardrail checks run every 15 minutes
\item \textbf{Alerting:} CRITICAL alerts (MAE, ROI, Win Rate) trigger PagerDuty + auto-stop; MEDIUM alerts (Calibration, Latency) send Slack notification
\item \textbf{Data Retention:} Full prediction logs retained for 6 months, aggregated metrics for 2 years (regulatory compliance)
\item \textbf{Access Control:} Dashboard read-only for analysts, alert configuration restricted to ops team
\end{tablenotes}
\end{table}

\paragraph{Sample Query: MAE by Group.}
The primary monitoring query computes rolling MAE for each group:

\begin{lstlisting}[language=SQL, style=codestyle, caption={Grafana query for real-time MAE monitoring}]
SELECT
  ab_group,
  DATE_TRUNC('hour', predicted_at) as time_bucket,
  AVG(ABS(prediction_mean - actual_value)) as mae,
  COUNT(*) as n_predictions
FROM ml.predictions
WHERE test_name = 'ensemble_v5_vs_v6_xgboost_ridge'
  AND actual_value IS NOT NULL
  AND predicted_at > NOW() - INTERVAL '24 hours'
GROUP BY ab_group, time_bucket
ORDER BY time_bucket ASC;
\end{lstlisting}

This query enables visualization of MAE trends over time, early detection of degradation, and validation of expected 7-yard improvement.

\subsection{Lessons Learned \& Best Practices}

Implementation of the November 2025 A/B test revealed several critical insights:

\paragraph{1. Deterministic Assignment is Non-Negotiable for Financial Applications.}
Randomized assignment (common in web A/B testing) violates reproducibility requirements for financial auditing. Cryptographic hashing (MD5) ensures:
\begin{itemize}[itemsep=2pt,parsep=0pt]
    \item Independent auditors can verify which model produced which bet
    \item Disputes can be resolved by re-running hash computation
    \item Regulatory compliance (betting license requires audit trail)
\end{itemize}

\paragraph{2. Guardrails Must Be Conservative to Prevent Catastrophic Loss.}
Initial guardrail design allowed 20\% MAE degradation, which would have cost \$4,000--\$10,000 over 4 weeks. Final design (10\% degradation, auto-stop within 60 seconds) limits max loss to \$2,000 and preserves bankroll for recovery.

\paragraph{3. Sequential Testing Saves Time Without Sacrificing Rigor.}
O'Brien-Fleming boundaries enable early stopping when treatment is clearly superior (Week 2 early stop possible) while maintaining family-wise Type I error rate at $\alpha = 0.05$. For the XGBoost v4 test, expected early stop at Week 2 saves 2 weeks of conservative baseline bets (\$1,000--\$2,000 opportunity cost).

\paragraph{4. Multi-Metric Success Criteria Prevent Single-Metric Optimization.}
MAE-only criteria would miss calibration breakdown (interval width collapse) or latency degradation (SLA violation). Six-metric framework (Table~\ref{tab:ab-success-criteria}) ensures production readiness across accuracy, business value, and operational constraints.

\paragraph{5. Real-Time Monitoring is Essential for Fast Intervention.}
15-minute guardrail checks (vs daily batches) enable 60-second rollback from degradation detection, preventing multi-hour losses. Grafana + PagerDuty integration automates alert → rollback workflow without human intervention.

\subsection{Statistical Power Analysis: Sample Size Determination}

The choice of $n = 100$ games per group balances statistical rigor, time-to-decision, and operational cost.

\paragraph{Power Calculation.}
For paired Diebold-Mariano test detecting 5-yard MAE improvement:

\begin{equation}
\begin{aligned}
\text{Power} &= \Phi \left( \frac{\mu_{\Delta} \sqrt{n}}{\sigma_{\Delta}} - z_{\alpha/2} \right) \\
&= \Phi \left( \frac{5.0 \sqrt{100}}{12.0} - 1.96 \right) \\
&= \Phi(4.17 - 1.96) = \Phi(2.21) = 0.987 \approx 98.7\%
\end{aligned}
\label{eq:power-calculation}
\end{equation}

where $\mu_{\Delta} = 5.0$ yards (minimum detectable effect), $\sigma_{\Delta} = 12.0$ yards (standard deviation of pairwise differences from 2022--2024 backtest), $n = 100$ games per group, and $z_{\alpha/2} = 1.96$ for $\alpha = 0.05$ two-tailed test.

This provides 98.7\% power (far exceeding 80\% standard), ensuring high probability of detecting 5-yard improvement if it truly exists.

\paragraph{Time-to-Decision Trade-Off.}
NFL regular season has $\sim$16--17 games per week, $\sim$20 active QBs per week, yielding $\sim$25--30 player-game predictions per week:

\begin{itemize}[itemsep=2pt,parsep=0pt]
    \item $n = 50$ per group: 2 weeks, 60\% power (acceptable for early stop with $\Delta\text{MAE} > 7$)
    \item $n = 100$ per group: 4 weeks, 98.7\% power (final analysis)
    \item $n = 150$ per group: 6 weeks, 99.9\% power (overkill, opportunity cost too high)
\end{itemize}

Four weeks balances rigor (98.7\% power) with time-to-market (mid-season deployment).

\subsection{Alternative Designs Considered}

Several alternative A/B test frameworks were evaluated and rejected:

\paragraph{Bandit Algorithms (Thompson Sampling, UCB).}
Multi-armed bandit algorithms dynamically allocate more traffic to superior model, minimizing regret (opportunity cost). Rejected because:
\begin{itemize}[itemsep=2pt,parsep=0pt]
    \item Requires probabilistic assignment (violates determinism for auditing)
    \item Complicates statistical inference (non-uniform treatment probabilities)
    \item Optimizes for short-term regret, not long-term deployment decision
    \item Bandit exploration can waste bets on inferior model early in test
\end{itemize}

\paragraph{Bayesian A/B Testing (Beta-Binomial Posterior).}
Bayesian framework models win rate as $\theta \sim \text{Beta}(\alpha, \beta)$, updates posterior after each game, computes $P(\theta_{\text{treatment}} > \theta_{\text{control}})$. Rejected because:
\begin{itemize}[itemsep=2pt,parsep=0pt]
    \item Requires prior elicitation (subjective, contentious)
    \item No universally accepted Bayesian analog to Diebold-Mariano test for forecast comparison
    \item Regulatory compliance easier with frequentist $p$-values
    \item Complexity not justified for 4-week test (Bayesian benefits emerge in continuous monitoring)
\end{itemize}

\paragraph{Interleaved Testing (Rank Fusion).}
Present user with predictions from both models, infer preference from click behavior. Not applicable to automated betting (no human in loop).

\subsection{Regulatory \& Ethical Considerations}

A/B testing with real money requires careful attention to compliance and ethics:

\paragraph{1. Informed Consent \& Transparency.}
For automated betting (no human bettors), transparency requirements differ from consumer-facing A/B tests:
\begin{itemize}[itemsep=2pt,parsep=0pt]
    \item Internal stakeholders (research team, operations) aware of test design
    \item Test plan documented and approved by risk committee
    \item Guardrails and auto-stop mechanisms disclosed to regulators
    \item No deception: both models are production-quality candidates
\end{itemize}

\paragraph{2. Data Privacy \& Security.}
All predictions and results logged to secure database with access controls:
\begin{itemize}[itemsep=2pt,parsep=0pt]
    \item Player IDs pseudonymized (do not expose NFL player personally identifiable information)
    \item Database credentials in environment variables (not hardcoded)
    \item Audit trail immutable (predictions cannot be modified post-game)
    \item Retention policy: 6 months full logs, 2 years aggregated metrics
\end{itemize}

\paragraph{3. Financial Responsibility: Kelly Criterion Compliance.}
Both control and treatment models must comply with Kelly criterion for bet sizing:
\begin{itemize}[itemsep=2pt,parsep=0pt]
    \item Maximum bet size: $f^* = \frac{p(b+1) - 1}{b}$ where $p$ is win probability, $b$ is odds
    \item Calibration requirement: 90\% intervals must have 90\% coverage (enables accurate $p$ estimation)
    \item Guardrail: Auto-stop if calibration drifts outside [85\%, 95\%]
    \item Bankroll protection: Max 2\% of bankroll per bet (conservative multiplier on Kelly)
\end{itemize}

Miscalibrated predictions violate Kelly optimality and risk ruin; hence calibration is a CRITICAL guardrail (Table~\ref{tab:ab-guardrails}).

\subsection{Expected November 2025 Test Outcome}

Based on XGBoost v4 multi-season validation (Section~\ref{sec:ensemble-evolution}), we anticipate:

\paragraph{Week 2 Early Stop (Most Likely Scenario).}
XGBoost v4 demonstrated 4.70 $\pm$ 0.40 MAE across 2022--2024 (88.31\% improvement, all $p < 0.000001$). This suggests:
\begin{itemize}[itemsep=2pt,parsep=0pt]
    \item Expected $\Delta\text{MAE} \approx 7$ yards (40 yards v5 $-$ 33 yards v6)
    \item $p$-value likely $< 0.001$ at $n = 50$ (exceeds Week 2 threshold $p < 0.01$)
    \item Early stop at Week 2 most probable outcome
    \item Saves 2 weeks opportunity cost (\$1,000--\$2,000)
\end{itemize}

\paragraph{Alternative Outcome: Week 4 Final Decision.}
If 2025 season exhibits higher variance (more upsets, defensive schemes adapting), effect size may diminish:
\begin{itemize}[itemsep=2pt,parsep=0pt]
    \item $\Delta\text{MAE} = 5$ yards (minimum threshold)
    \item $p = 0.03$ at $n = 100$ (passes final threshold $p < 0.034$)
    \item Deploy at Week 4 after collecting 100 games
\end{itemize}

\paragraph{Unlikely Outcome: Guardrail Violation.}
If XGBoost v4 fails to generalize to 2025 (e.g., due to rule changes, COVID-era data contamination):
\begin{itemize}[itemsep=2pt,parsep=0pt]
    \item MAE degradation $> 10\%$ triggers auto-stop
    \item Rollback to v5 within 60 seconds
    \item Post-mortem analysis: investigate data drift, re-validate training pipeline
\end{itemize}

Historical precedent (3 years of consistent performance) makes guardrail violation low probability ($< 5\%$).

\subsection{Summary: Principled Model Validation for Production Deployment}

The A/B testing framework instantiates key principles of rigorous machine learning engineering:

\begin{enumerate}[itemsep=3pt,parsep=0pt]
    \item \textbf{Determinism:} Cryptographic hashing ensures reproducible assignment and audit compliance
    \item \textbf{Statistical Rigor:} Diebold-Mariano test, O'Brien-Fleming sequential boundaries, Bonferroni correction control Type I error
    \item \textbf{Risk Management:} Conservative guardrails (10\% MAE degradation) with 60-second auto-stop prevent catastrophic loss
    \item \textbf{Multi-Metric Validation:} Six success criteria (MAE, RMSE, Win Rate, ROI, Calibration, Latency) ensure production readiness
    \item \textbf{Operational Excellence:} Real-time Grafana dashboards, PagerDuty alerts, automated rollback enable fast intervention
    \item \textbf{Efficiency:} Sequential testing with early stop opportunities balances rigor with time-to-market
\end{enumerate}

This infrastructure generalizes beyond the November 2025 XGBoost v4 test to any ensemble upgrade: future temperature scaling refinements (Section~\ref{sec:temperature-scaling}), state-space model improvements, or hybrid architecture experiments can leverage the same framework, accumulating institutional knowledge for confident production deployment decisions.

The broader lesson: \emph{production machine learning requires engineering rigor beyond algorithmic innovation}. While XGBoost v4's 88.31\% improvement is algorithmically impressive, the A/B testing framework---deterministic routing, sequential testing, guardrails, monitoring---converts research findings into production value safely and reproducibly. This infrastructure investment pays compounding dividends across all future model iterations.
